\chapter*{Глоссарий}
\addcontentsline{toc}{chapter}{Глоссарий}
\markboth{Глоссарий}{}

\newcommand{\gloss}[2]{\noindent\textbf{#1}~--- #2\par\vspace{3pt}}

\section*{Платформа и язык}

\gloss{Байт-код}{промежуточное представление программы, получаемое компилятором
\texttt{javac} из исходного текста. Не привязан к процессору и операционной
системе, исполняется виртуальной машиной.}

\gloss{Виртуальная машина Java (JVM)}{программа, исполняющая байт-код. Отвечает
за загрузку классов, управление памятью, сборку мусора и оптимизацию во время
выполнения.}

\gloss{Загрузчик классов (\eng{ClassLoader})}{компонент JVM, находящий и
подгружающий классы в память по мере обращения к ним. Работает по модели
делегирования: сначала спрашивает родительский загрузчик.}

\gloss{JIT-компиляция}{перевод часто исполняемых участков байт-кода в машинный
код непосредственно во время работы программы. Объясняет, почему Java-приложение
разгоняется через несколько секунд после запуска.}

\gloss{Сборщик мусора (\eng{Garbage Collector})}{механизм автоматического
освобождения памяти от объектов, на которые не осталось ссылок.}

\gloss{JDK}{комплект разработчика: компилятор, стандартная библиотека, среда
выполнения и вспомогательные инструменты. JRE~--- только среда выполнения,
компилировать в ней нельзя.}

\gloss{Примитивный тип}{тип, значение которого хранится непосредственно
(\texttt{int}, \texttt{double}, \texttt{boolean} и другие). Ссылочный тип хранит
адрес объекта в куче.}

\section*{Объектно-ориентированное программирование}

\gloss{Класс}{описание структуры и поведения объектов: полей и методов. Объект~---
конкретный экземпляр класса, созданный оператором \texttt{new}.}

\gloss{Инкапсуляция}{сокрытие внутреннего устройства объекта за ограниченным
набором методов доступа. Позволяет менять реализацию, не ломая код, который
объектом пользуется.}

\gloss{Наследование}{механизм, при котором класс перенимает поля и методы
другого класса и может их дополнять или переопределять.}

\gloss{Полиморфизм}{способность вызвать метод через ссылку базового типа и
получить поведение фактического класса объекта.}

\gloss{Абстракция}{выделение существенных свойств объекта и отбрасывание
несущественных для конкретной задачи.}

\gloss{Композиция}{построение объекта из других объектов, хранимых в его полях.
В большинстве случаев предпочтительнее наследования: связь слабее, реализацию
можно подменить во время выполнения.}

\gloss{Интерфейс}{контракт, описывающий, какие методы обязан предоставить
реализующий его класс, без указания, как именно.}

\gloss{Sealed-класс}{класс, явно перечисляющий в \texttt{permits} те классы,
которым разрешено его наследовать. Даёт закрытую иерархию, известную компилятору.}

\section*{Ошибки и обобщения}

\gloss{Исключение}{объект, описывающий нештатную ситуацию. Проверяемые
(\eng{checked}) исключения компилятор требует обработать или объявить;
непроверяемые (наследники \texttt{RuntimeException}) --- нет.}

\gloss{\texttt{try-with-resources}}{конструкция, автоматически закрывающая
ресурсы, реализующие \texttt{AutoCloseable}, при выходе из блока~--- в том числе
при исключении.}

\gloss{Обобщения (\eng{Generics})}{параметризация классов и методов типами.
Переносят часть ошибок с этапа выполнения на этап компиляции.}

\gloss{Стирание типов (\eng{type erasure})}{удаление информации о параметрах
типа при компиляции. Причина того, что во время выполнения \texttt{List<String>}
и \texttt{List<Integer>}~--- один и тот же тип.}

\section*{Коллекции, ввод-вывод, многопоточность}

\gloss{Коллекция}{объект, хранящий группу элементов. \texttt{List}~--- упорядоченный
список с дубликатами, \texttt{Set}~--- множество без дубликатов,
\texttt{Map}~--- набор пар «ключ~--- значение».}

\gloss{Поток ввода-вывода}{канал передачи данных. Байтовые потоки
(\texttt{InputStream}, \texttt{OutputStream}) работают с двоичными данными,
символьные (\texttt{Reader}, \texttt{Writer})~--- с текстом с учётом кодировки.}

\gloss{Поток выполнения (\eng{Thread})}{независимая последовательность команд
внутри процесса. Потоки одного процесса делят общую кучу, но каждый имеет
собственный стек.}

\gloss{Синхронизация}{механизмы упорядочивания доступа нескольких потоков к
общим данным: \texttt{synchronized}, блокировки, атомарные типы.}

\gloss{Stream API}{средства декларативной обработки данных конвейером из
промежуточных и терминальных операций. Не имеет отношения ни к потокам
ввода-вывода, ни к потокам выполнения~--- совпадает только английское слово.}

\section*{Базы данных}

\gloss{JDBC}{стандартный интерфейс доступа к реляционным базам данных из Java.}

\gloss{\texttt{PreparedStatement}}{подготовленный запрос с параметрами.
Основное средство защиты от SQL-инъекций: параметры передаются отдельно от
текста запроса.}

\gloss{ORM}{объектно-реляционное отображение~--- автоматическое сопоставление
классов таблицам, а полей~--- столбцам.}

\gloss{Транзакция}{группа операций, выполняемых как единое целое: либо все,
либо ни одной. Свойства~--- атомарность, согласованность, изолированность,
долговечность (ACID).}

\gloss{Уровень изоляции}{степень защищённости транзакции от параллельных
изменений. Чем выше уровень, тем меньше аномалий и тем ниже параллелизм.}

\section*{Spring и веб}

\gloss{Инверсия управления (IoC)}{принцип, при котором созданием объектов и
связыванием их между собой занимается контейнер, а не сам код.}

\gloss{Внедрение зависимостей (DI)}{способ реализации IoC: объект получает
зависимости извне~--- через конструктор, сеттер или поле~--- вместо того чтобы
создавать их самому.}

\gloss{Бин (\eng{bean})}{объект, жизненным циклом которого управляет
контейнер Spring.}

\gloss{HTTP}{текстовый протокол запрос-ответ, лежащий в основе веба. Не хранит
состояние между запросами.}

\gloss{REST}{архитектурный стиль, в котором данные представлены ресурсами с
собственными адресами, а действия над ними выражаются методами HTTP.}

\gloss{DTO}{объект передачи данных~--- отдельный от сущности класс для обмена
с внешним миром. Разрывает связь между внутренней моделью и внешним API.}

\gloss{Сервлет}{Java-класс, обрабатывающий HTTP-запросы внутри контейнера
приложений. Spring MVC построен поверх сервлетов.}

\section*{Качество кода}

\gloss{Модульный тест}{автоматическая проверка отдельного участка кода в
изоляции от остальной системы.}

\gloss{Мок (\eng{mock})}{подставной объект, заменяющий настоящую зависимость
в тесте и позволяющий проверить, как с ней взаимодействовали.}

\gloss{Javadoc}{формат документирующих комментариев и утилита, собирающая из
них HTML-документацию.}

\gloss{Паттерн проектирования}{типовое решение часто встречающейся задачи
проектирования, описанное так, чтобы его можно было применить повторно.}

\gloss{Антипаттерн}{распространённое решение, которое выглядит удачным, но
приводит к предсказуемым проблемам.}

\gloss{Рефакторинг}{изменение внутренней структуры кода без изменения его
внешнего поведения.}

\section*{Контроль версий и процессы}

\gloss{Коммит}{зафиксированное состояние проекта вместе с указателем на
предыдущий коммит. История проекта~--- граф таких состояний.}

\gloss{Ветка}{подвижный указатель на коммит. Позволяет вести несколько линий
разработки параллельно.}

\gloss{Слияние (\eng{merge})}{объединение двух линий разработки. При
несовместимых изменениях возникает конфликт, который разрешается вручную.}

\gloss{CI/CD}{автоматическая сборка и проверка кода при каждом изменении
(непрерывная интеграция) и автоматическая доставка проверенной версии
(непрерывная поставка).}

\gloss{Scrum}{способ организации работы короткими итерациями~--- спринтами~---
с фиксированным набором ролей, артефактов и событий.}
